% TOP_PROBLEM: {"id": 7200050, "title": "The Kobon triangle problem on triangles in line arrangements", "category_id": 6, "category": {"id": 6, "name": "geometry", "display_name": "Geometry", "description": "Euclidean and non-Euclidean geometry, geometric structures.", "slug": "geometry", "order_index": 6, "created_at": "2026-07-31T15:26:25.671Z"}, "status": "open"}
% FIRST_POSTED: null
% ENTRY_KIND: statement_only
% Imported from: batch19.tex; UnsolvedMath ID 7200050
\documentclass[11pt]{article}
\usepackage{fontspec}
\setmainfont{Latin Modern Roman}
\setsansfont{Latin Modern Sans}
\usepackage{amsmath,amssymb,mathtools}
\usepackage{unicode-math}
\setmathfont{Latin Modern Math}
\usepackage{xurl}
\usepackage[hidelinks]{hyperref}
\newcommand{\sref}[2]{\hyperlink{source:#1:#2}{[S#2]}}
\newcommand{\cataloguescope}[1]{\par\noindent\textbf{Question.} #1\par}
\begin{document}
% UnsolvedMath ID 7200050; source catalogue code AMR-071-0050
\setcounter{section}{7200049}
\section{The Kobon triangle problem on triangles in line arrangements}\label{7200050}
{\small\sffamily UnsolvedMath ID 7200050 · Geometry}
\subsection{Definitions and mathematical statement}

\paragraph{Shared notation.}$\mathbb N=\{1,2,\ldots\}$, and $\mathbb Z,\mathbb Q,\mathbb R,\mathbb C$ denote the integers, rationals, reals and complex numbers. Any other conventions are specified locally.


An arrangement consists of $n$ distinct straight lines in $\mathbb R^2$. The lines subdivide the complement into regions. A triangular region is a bounded region whose closure has three sides, each lying on a line of the arrangement; its interior is crossed by no arrangement line. Distinct such regions have disjoint interiors. Parallel lines and multiple concurrent lines are allowed; no restriction to general position is imposed. Let $T(n)$ be the largest possible number of triangular regions over all such arrangements.
\paragraph{Statement recorded in the supplied source.}
Determine $T(n)$, the maximum number of nonoverlapping triangles that can be formed by an arrangement of $n$ straight lines, for each $n\geq1$. \sref{7200050}{2}
\subsection{Short English statement}
How many triangular cells can an arrangement of a prescribed number of straight lines produce?
\subsection{Sources}
\begin{description}
\item[\hypertarget{source:7200050:1}{S1}] ULAM.AI and the UnsolvedMath Contributors, UnsolvedMath: A Curated Collection of Open Mathematics Problems, record 7200050 (AMR-071-0050). CC BY 4.0. \url{https://huggingface.co/datasets/ulamai/UnsolvedMath}.
\item[\hypertarget{source:7200050:2}{S2}] Pavlo Savchuk, Constructing Optimal Kobon Triangle Arrangements via Table Encoding, SAT Solving, and Heuristic Straightening (2026 version of 2025 preprint). Sections 1 and 2.1–2.2, pp. 1–3. \url{https://arxiv.org/pdf/2507.07951}.
\end{description}
\label{end:7200050}
\end{document}
